
\documentclass{article}
\usepackage{colortbl}
\usepackage{makecell}
\usepackage{multirow}
\usepackage{supertabular}

\begin{document}

\newcounter{utterance}

\centering \large Interaction Transcript for game `matchit\_ascii', experiment `different\_grid', episode 6 with Qwen3.6{-}35B{-}A3B{-}FP8{-}without{-}reasoning.
\vspace{24pt}

{ \footnotesize  \setcounter{utterance}{1}
\setlength{\tabcolsep}{0pt}
\begin{supertabular}{c@{$\;$}|p{.15\linewidth}@{}p{.15\linewidth}p{.15\linewidth}p{.15\linewidth}p{.15\linewidth}p{.15\linewidth}}
    \# & $\;$A & \multicolumn{4}{c}{Game Master} & $\;\:$B\\
    \hline

    \theutterance \stepcounter{utterance}  
    & & \multicolumn{4}{p{0.6\linewidth}}{
        \cellcolor[rgb]{0.9,0.9,0.9}{
            \makecell[{{p{\linewidth}}}]{
                \texttt{\tiny{[P1$\langle$GM]}}
                \texttt{Jūs piedalāties kopīgā uzminēšanas spēlē.} \\
\texttt{Mērķis ir noskaidrot, vai šis režģis un cits attēls, ko redzu tikai es, ir vienādi.} \\
\texttt{Lūdzu, vispirms aprakstiet savu režģi. Pēc tam es sniegšu savu aprakstu, un mēs varēsim viens otram uzdot jautājumus par režģiem, lai noskaidrotu, vai tie ir vienādi. Tagad sāciet savu īso režģa aprakstu ar "APRAKSTS:", kam seko apraksts. Nepievienojiet neko citu. Šis ir jūsu režģis:} \\
\texttt{X ▢ ▢ ▢ X} \\
\texttt{X ▢ ▢ ▢ X} \\
\texttt{X X X X X} \\
\texttt{X ▢ ▢ ▢ X} \\
\texttt{X ▢ ▢ ▢ X} \\
            }
        }
    }
    & \\ \\

    \theutterance \stepcounter{utterance}  
    & \multicolumn{4}{p{0.6\linewidth}}{
        \cellcolor[rgb]{0.9,0.9,0.9}{
            \makecell[{{p{\linewidth}}}]{
                \texttt{\tiny{[P1$\rangle$GM]}}
                \texttt{APRAKSTS: 5x5 režģis, kurā trešā rinda un trešais stabs ir pilnībā aizpildīti ar X, bet pārējie laukumi ir tukši (▢).} \\
            }
        }
    }
    & & \\ \\

    \theutterance \stepcounter{utterance}  
    & & & \multicolumn{2}{p{0.3\linewidth}}{
        \cellcolor[rgb]{0.9,0.9,0.9}{
            \makecell[{{p{\linewidth}}}]{
                \texttt{\tiny{[GM$|$GM]}}
                \texttt{continue} \\
            }
        }
    }
    & & \\ \\

    \theutterance \stepcounter{utterance}  
    & & \multicolumn{4}{p{0.6\linewidth}}{
        \cellcolor[rgb]{0.9,0.9,0.9}{
            \makecell[{{p{\linewidth}}}]{
                \texttt{\tiny{[GM$\rangle$P2]}}
                \texttt{Jūs piedalāties kopīgā uzminēšanas spēlē.} \\
\texttt{Mērķis ir noskaidrot, vai šis režģis un cits attēls, ko redzu tikai es, ir vienādi.} \\
\texttt{Lūdzu, vispirms aprakstiet savu režģi. Pēc tam es sniegšu savu aprakstu, un mēs varēsim viens otram uzdot jautājumus par režģiem, lai noskaidrotu, vai tie ir vienādi. Tagad sāciet savu īso režģa aprakstu ar "APRAKSTS:", kam seko apraksts. Nepievienojiet neko citu. Šis ir jūsu režģis:} \\
\texttt{▢ ▢ ▢ ▢ ▢} \\
\texttt{X X X ▢ X} \\
\texttt{X X ▢ X X} \\
\texttt{X ▢ X X X} \\
\texttt{▢ ▢ ▢ ▢ ▢} \\
            }
        }
    }
    & \\ \\

    \theutterance \stepcounter{utterance}  
    & & & \multicolumn{4}{p{0.6\linewidth}}{
        \cellcolor[rgb]{0.9,0.9,0.9}{
            \makecell[{{p{\linewidth}}}]{
                \texttt{\tiny{[GM$\langle$P2]}}
                \texttt{APRAKSTS: 5x5 režģis ar tukšām šūnām (▢) augšējā un apakšējā rindā. Vidējās trīs rindas satur krustenotus simbolus (X) un tukšas šūnas (▢) noteiktā secībā, veidojot simetrisku rakstu attiecībā pret centrālo kolonnu.} \\
            }
        }
    }
     \\ \\

    \theutterance \stepcounter{utterance}  
    & & & \multicolumn{2}{p{0.3\linewidth}}{
        \cellcolor[rgb]{0.9,0.9,0.9}{
            \makecell[{{p{\linewidth}}}]{
                \texttt{\tiny{[GM$|$GM]}}
                \texttt{continue} \\
            }
        }
    }
    & & \\ \\

    \theutterance \stepcounter{utterance}  
    & & \multicolumn{4}{p{0.6\linewidth}}{
        \cellcolor[rgb]{0.9,0.9,0.9}{
            \makecell[{{p{\linewidth}}}]{
                \texttt{\tiny{[GM$\rangle$P2]}}
                \texttt{Šis ir mans APRAKSTSAPRAKSTS: 5x5 režģis, kurā trešā rinda un trešais stabs ir pilnībā aizpildīti ar X, bet pārējie laukumi ir tukši (▢).} \\
\texttt{Tagad uzdodiet jautājumu, lai noskaidrotu jaunus aspektus par manu režģi, kas var atšķirties no jūsu režģa. Sāciet ar "JAUTĀJUMS:" un nepievienojiet neko citu.} \\
            }
        }
    }
    & \\ \\

    \theutterance \stepcounter{utterance}  
    & & & \multicolumn{4}{p{0.6\linewidth}}{
        \cellcolor[rgb]{0.9,0.9,0.9}{
            \makecell[{{p{\linewidth}}}]{
                \texttt{\tiny{[GM$\langle$P2]}}
                \texttt{JAUTĀJUMS: Vai jūsu režģī ir X arī pirmajā un piektajā rindā, vai arī tās ir pilnībā tukšas?} \\
            }
        }
    }
     \\ \\

    \theutterance \stepcounter{utterance}  
    & & & \multicolumn{2}{p{0.3\linewidth}}{
        \cellcolor[rgb]{0.9,0.9,0.9}{
            \makecell[{{p{\linewidth}}}]{
                \texttt{\tiny{[GM$|$GM]}}
                \texttt{continue} \\
            }
        }
    }
    & & \\ \\

    \theutterance \stepcounter{utterance}  
    & & \multicolumn{4}{p{0.6\linewidth}}{
        \cellcolor[rgb]{0.9,0.9,0.9}{
            \makecell[{{p{\linewidth}}}]{
                \texttt{\tiny{[P1$\langle$GM]}}
                \texttt{Šis ir mans APRAKSTSAPRAKSTS: 5x5 režģis ar tukšām šūnām (▢) augšējā un apakšējā rindā. Vidējās trīs rindas satur krustenotus simbolus (X) un tukšas šūnas (▢) noteiktā secībā, veidojot simetrisku rakstu attiecībā pret centrālo kolonnu.} \\
\texttt{JAUTĀJUMS: Vai jūsu režģī ir X arī pirmajā un piektajā rindā, vai arī tās ir pilnībā tukšas?Sāciet savu atbildi ar "ATBILDE:" un nepievienojiet neko citu.} \\
            }
        }
    }
    & \\ \\

    \theutterance \stepcounter{utterance}  
    & \multicolumn{4}{p{0.6\linewidth}}{
        \cellcolor[rgb]{0.9,0.9,0.9}{
            \makecell[{{p{\linewidth}}}]{
                \texttt{\tiny{[P1$\rangle$GM]}}
                \texttt{ATBILDE: Jā, pirmajā un piektajā rindā ir X pirmajā un piektajā kolonnā, bet pārējie laukumi ir tukši.} \\
            }
        }
    }
    & & \\ \\

    \theutterance \stepcounter{utterance}  
    & & & \multicolumn{2}{p{0.3\linewidth}}{
        \cellcolor[rgb]{0.9,0.9,0.9}{
            \makecell[{{p{\linewidth}}}]{
                \texttt{\tiny{[GM$|$GM]}}
                \texttt{continue} \\
            }
        }
    }
    & & \\ \\

    \theutterance \stepcounter{utterance}  
    & & \multicolumn{4}{p{0.6\linewidth}}{
        \cellcolor[rgb]{0.9,0.9,0.9}{
            \makecell[{{p{\linewidth}}}]{
                \texttt{\tiny{[P1$\langle$GM]}}
                \texttt{Tagad uzdodiet jautājumu, lai noskaidrotu jaunus aspektus par manu režģi, kas var atšķirties no jūsu režģa. Sāciet ar "JAUTĀJUMS:" un nepievienojiet neko citu.} \\
            }
        }
    }
    & \\ \\

    \theutterance \stepcounter{utterance}  
    & \multicolumn{4}{p{0.6\linewidth}}{
        \cellcolor[rgb]{0.9,0.9,0.9}{
            \makecell[{{p{\linewidth}}}]{
                \texttt{\tiny{[P1$\rangle$GM]}}
                \texttt{JAUTĀJUMS: Vai jūsu režģī X ir arī otrajā un ceturtajā rindā, vai arī tās satur tikai tukšas šūnas?} \\
            }
        }
    }
    & & \\ \\

    \theutterance \stepcounter{utterance}  
    & & & \multicolumn{2}{p{0.3\linewidth}}{
        \cellcolor[rgb]{0.9,0.9,0.9}{
            \makecell[{{p{\linewidth}}}]{
                \texttt{\tiny{[GM$|$GM]}}
                \texttt{continue} \\
            }
        }
    }
    & & \\ \\

    \theutterance \stepcounter{utterance}  
    & & \multicolumn{4}{p{0.6\linewidth}}{
        \cellcolor[rgb]{0.9,0.9,0.9}{
            \makecell[{{p{\linewidth}}}]{
                \texttt{\tiny{[GM$\rangle$P2]}}
                \texttt{ATBILDE: Jā, pirmajā un piektajā rindā ir X pirmajā un piektajā kolonnā, bet pārējie laukumi ir tukši.} \\
\texttt{JAUTĀJUMS: Vai jūsu režģī X ir arī otrajā un ceturtajā rindā, vai arī tās satur tikai tukšas šūnas?Sāciet savu atbildi ar "ATBILDE:" un nepievienojiet neko citu.} \\
            }
        }
    }
    & \\ \\

    \theutterance \stepcounter{utterance}  
    & & & \multicolumn{4}{p{0.6\linewidth}}{
        \cellcolor[rgb]{0.9,0.9,0.9}{
            \makecell[{{p{\linewidth}}}]{
                \texttt{\tiny{[GM$\langle$P2]}}
                \texttt{ATBILDE: Nē, otrajā un ceturtajā rindā nav X pirmajā un piektajā kolonnā; tās ir tukšas (▢).} \\
            }
        }
    }
     \\ \\

    \theutterance \stepcounter{utterance}  
    & & & \multicolumn{2}{p{0.3\linewidth}}{
        \cellcolor[rgb]{0.9,0.9,0.9}{
            \makecell[{{p{\linewidth}}}]{
                \texttt{\tiny{[GM$|$GM]}}
                \texttt{continue} \\
            }
        }
    }
    & & \\ \\

    \theutterance \stepcounter{utterance}  
    & & \multicolumn{4}{p{0.6\linewidth}}{
        \cellcolor[rgb]{0.9,0.9,0.9}{
            \makecell[{{p{\linewidth}}}]{
                \texttt{\tiny{[GM$\rangle$P2]}}
                \texttt{Tagad uzdodiet jautājumu, lai noskaidrotu jaunus aspektus par manu režģi, kas var atšķirties no jūsu režģa. Sāciet ar "JAUTĀJUMS:" un nepievienojiet neko citu.} \\
            }
        }
    }
    & \\ \\

    \theutterance \stepcounter{utterance}  
    & & & \multicolumn{4}{p{0.6\linewidth}}{
        \cellcolor[rgb]{0.9,0.9,0.9}{
            \makecell[{{p{\linewidth}}}]{
                \texttt{\tiny{[GM$\langle$P2]}}
                \texttt{JAUTĀJUMS: Vai jūsu režģī otrajā un ceturtajā rindā X atrodas tikai otrajā un ceturtajā kolonnā, vai arī ir vēl kādi X šajās rindās?} \\
            }
        }
    }
     \\ \\

    \theutterance \stepcounter{utterance}  
    & & & \multicolumn{2}{p{0.3\linewidth}}{
        \cellcolor[rgb]{0.9,0.9,0.9}{
            \makecell[{{p{\linewidth}}}]{
                \texttt{\tiny{[GM$|$GM]}}
                \texttt{continue} \\
            }
        }
    }
    & & \\ \\

    \theutterance \stepcounter{utterance}  
    & & \multicolumn{4}{p{0.6\linewidth}}{
        \cellcolor[rgb]{0.9,0.9,0.9}{
            \makecell[{{p{\linewidth}}}]{
                \texttt{\tiny{[P1$\langle$GM]}}
                \texttt{ATBILDE: Nē, otrajā un ceturtajā rindā nav X pirmajā un piektajā kolonnā; tās ir tukšas (▢).} \\
\texttt{JAUTĀJUMS: Vai jūsu režģī otrajā un ceturtajā rindā X atrodas tikai otrajā un ceturtajā kolonnā, vai arī ir vēl kādi X šajās rindās?Sāciet savu atbildi ar "ATBILDE:" un nepievienojiet neko citu.} \\
            }
        }
    }
    & \\ \\

    \theutterance \stepcounter{utterance}  
    & \multicolumn{4}{p{0.6\linewidth}}{
        \cellcolor[rgb]{0.9,0.9,0.9}{
            \makecell[{{p{\linewidth}}}]{
                \texttt{\tiny{[P1$\rangle$GM]}}
                \texttt{ATBILDE: Nē, otrajā un ceturtajā rindā nav X otrajā un ceturtajā kolonnā; tās ir pilnībā tukšas (▢).} \\
            }
        }
    }
    & & \\ \\

    \theutterance \stepcounter{utterance}  
    & & & \multicolumn{2}{p{0.3\linewidth}}{
        \cellcolor[rgb]{0.9,0.9,0.9}{
            \makecell[{{p{\linewidth}}}]{
                \texttt{\tiny{[GM$|$GM]}}
                \texttt{continue} \\
            }
        }
    }
    & & \\ \\

    \theutterance \stepcounter{utterance}  
    & & \multicolumn{4}{p{0.6\linewidth}}{
        \cellcolor[rgb]{0.9,0.9,0.9}{
            \makecell[{{p{\linewidth}}}]{
                \texttt{\tiny{[P1$\langle$GM]}}
                \texttt{Tagad uzdodiet jautājumu, lai noskaidrotu jaunus aspektus par manu režģi, kas var atšķirties no jūsu režģa. Sāciet ar "JAUTĀJUMS:" un nepievienojiet neko citu.} \\
            }
        }
    }
    & \\ \\

    \theutterance \stepcounter{utterance}  
    & \multicolumn{4}{p{0.6\linewidth}}{
        \cellcolor[rgb]{0.9,0.9,0.9}{
            \makecell[{{p{\linewidth}}}]{
                \texttt{\tiny{[P1$\rangle$GM]}}
                \texttt{JAUTĀJUMS: Vai jūsu režģī centrālajā kolonnā (trešajā kolonnā) ir X visās rindās, vai arī kādā no rindām šī kolonna ir tukša?} \\
            }
        }
    }
    & & \\ \\

    \theutterance \stepcounter{utterance}  
    & & & \multicolumn{2}{p{0.3\linewidth}}{
        \cellcolor[rgb]{0.9,0.9,0.9}{
            \makecell[{{p{\linewidth}}}]{
                \texttt{\tiny{[GM$|$GM]}}
                \texttt{continue} \\
            }
        }
    }
    & & \\ \\

    \theutterance \stepcounter{utterance}  
    & & \multicolumn{4}{p{0.6\linewidth}}{
        \cellcolor[rgb]{0.9,0.9,0.9}{
            \makecell[{{p{\linewidth}}}]{
                \texttt{\tiny{[GM$\rangle$P2]}}
                \texttt{ATBILDE: Nē, otrajā un ceturtajā rindā nav X otrajā un ceturtajā kolonnā; tās ir pilnībā tukšas (▢).} \\
\texttt{JAUTĀJUMS: Vai jūsu režģī centrālajā kolonnā (trešajā kolonnā) ir X visās rindās, vai arī kādā no rindām šī kolonna ir tukša?Sāciet savu atbildi ar "ATBILDE:" un nepievienojiet neko citu.} \\
            }
        }
    }
    & \\ \\

    \theutterance \stepcounter{utterance}  
    & & & \multicolumn{4}{p{0.6\linewidth}}{
        \cellcolor[rgb]{0.9,0.9,0.9}{
            \makecell[{{p{\linewidth}}}]{
                \texttt{\tiny{[GM$\langle$P2]}}
                \texttt{ATBILDE: Jā, centrālajā kolonnā (trešajā kolonnā) ir X visās rindās, ieskaitot pirmo un pēdējo rindu.} \\
            }
        }
    }
     \\ \\

    \theutterance \stepcounter{utterance}  
    & & & \multicolumn{2}{p{0.3\linewidth}}{
        \cellcolor[rgb]{0.9,0.9,0.9}{
            \makecell[{{p{\linewidth}}}]{
                \texttt{\tiny{[GM$|$GM]}}
                \texttt{continue} \\
            }
        }
    }
    & & \\ \\

    \theutterance \stepcounter{utterance}  
    & & \multicolumn{4}{p{0.6\linewidth}}{
        \cellcolor[rgb]{0.9,0.9,0.9}{
            \makecell[{{p{\linewidth}}}]{
                \texttt{\tiny{[GM$\rangle$P2]}}
                \texttt{Tagad uzdodiet jautājumu, lai noskaidrotu jaunus aspektus par manu režģi, kas var atšķirties no jūsu režģa. Sāciet ar "JAUTĀJUMS:" un nepievienojiet neko citu.} \\
            }
        }
    }
    & \\ \\

    \theutterance \stepcounter{utterance}  
    & & & \multicolumn{4}{p{0.6\linewidth}}{
        \cellcolor[rgb]{0.9,0.9,0.9}{
            \makecell[{{p{\linewidth}}}]{
                \texttt{\tiny{[GM$\langle$P2]}}
                \texttt{JAUTĀJUMS: Vai jūsu režģī otrajā un ceturtajā rindā ir X trešajā kolonnā?} \\
            }
        }
    }
     \\ \\

    \theutterance \stepcounter{utterance}  
    & & & \multicolumn{2}{p{0.3\linewidth}}{
        \cellcolor[rgb]{0.9,0.9,0.9}{
            \makecell[{{p{\linewidth}}}]{
                \texttt{\tiny{[GM$|$GM]}}
                \texttt{continue} \\
            }
        }
    }
    & & \\ \\

    \theutterance \stepcounter{utterance}  
    & & \multicolumn{4}{p{0.6\linewidth}}{
        \cellcolor[rgb]{0.9,0.9,0.9}{
            \makecell[{{p{\linewidth}}}]{
                \texttt{\tiny{[P1$\langle$GM]}}
                \texttt{ATBILDE: Jā, centrālajā kolonnā (trešajā kolonnā) ir X visās rindās, ieskaitot pirmo un pēdējo rindu.} \\
\texttt{JAUTĀJUMS: Vai jūsu režģī otrajā un ceturtajā rindā ir X trešajā kolonnā?Sāciet savu atbildi ar "ATBILDE:" un nepievienojiet neko citu.} \\
            }
        }
    }
    & \\ \\

    \theutterance \stepcounter{utterance}  
    & \multicolumn{4}{p{0.6\linewidth}}{
        \cellcolor[rgb]{0.9,0.9,0.9}{
            \makecell[{{p{\linewidth}}}]{
                \texttt{\tiny{[P1$\rangle$GM]}}
                \texttt{ATBILDE: Jā, otrajā un ceturtajā rindā X atrodas trešajā kolonnā.} \\
            }
        }
    }
    & & \\ \\

    \theutterance \stepcounter{utterance}  
    & & & \multicolumn{2}{p{0.3\linewidth}}{
        \cellcolor[rgb]{0.9,0.9,0.9}{
            \makecell[{{p{\linewidth}}}]{
                \texttt{\tiny{[GM$|$GM]}}
                \texttt{continue} \\
            }
        }
    }
    & & \\ \\

    \theutterance \stepcounter{utterance}  
    & & \multicolumn{4}{p{0.6\linewidth}}{
        \cellcolor[rgb]{0.9,0.9,0.9}{
            \makecell[{{p{\linewidth}}}]{
                \texttt{\tiny{[P1$\langle$GM]}}
                \texttt{Tagad uzdodiet jautājumu, lai noskaidrotu jaunus aspektus par manu režģi, kas var atšķirties no jūsu režģa. Sāciet ar "JAUTĀJUMS:" un nepievienojiet neko citu.} \\
            }
        }
    }
    & \\ \\

    \theutterance \stepcounter{utterance}  
    & \multicolumn{4}{p{0.6\linewidth}}{
        \cellcolor[rgb]{0.9,0.9,0.9}{
            \makecell[{{p{\linewidth}}}]{
                \texttt{\tiny{[P1$\rangle$GM]}}
                \texttt{JAUTĀJUMS: Vai jūsu režģī pirmajā un piektajā rindā X atrodas tikai pirmajā un piektajā kolonnā, vai arī ir vēl kādi X šajās rindās?} \\
            }
        }
    }
    & & \\ \\

    \theutterance \stepcounter{utterance}  
    & & & \multicolumn{2}{p{0.3\linewidth}}{
        \cellcolor[rgb]{0.9,0.9,0.9}{
            \makecell[{{p{\linewidth}}}]{
                \texttt{\tiny{[GM$|$GM]}}
                \texttt{continue} \\
            }
        }
    }
    & & \\ \\

    \theutterance \stepcounter{utterance}  
    & & \multicolumn{4}{p{0.6\linewidth}}{
        \cellcolor[rgb]{0.9,0.9,0.9}{
            \makecell[{{p{\linewidth}}}]{
                \texttt{\tiny{[GM$\rangle$P2]}}
                \texttt{ATBILDE: Jā, otrajā un ceturtajā rindā X atrodas trešajā kolonnā.} \\
\texttt{JAUTĀJUMS: Vai jūsu režģī pirmajā un piektajā rindā X atrodas tikai pirmajā un piektajā kolonnā, vai arī ir vēl kādi X šajās rindās?Sāciet savu atbildi ar "ATBILDE:" un nepievienojiet neko citu.} \\
            }
        }
    }
    & \\ \\

    \theutterance \stepcounter{utterance}  
    & & & \multicolumn{4}{p{0.6\linewidth}}{
        \cellcolor[rgb]{0.9,0.9,0.9}{
            \makecell[{{p{\linewidth}}}]{
                \texttt{\tiny{[GM$\langle$P2]}}
                \texttt{ATBILDE: Nē, pirmajā un piektajā rindā nav X pirmajā un piektajā kolonnā; tās ir tukšas (▢).} \\
            }
        }
    }
     \\ \\

    \theutterance \stepcounter{utterance}  
    & & & \multicolumn{2}{p{0.3\linewidth}}{
        \cellcolor[rgb]{0.9,0.9,0.9}{
            \makecell[{{p{\linewidth}}}]{
                \texttt{\tiny{[GM$|$GM]}}
                \texttt{continue} \\
            }
        }
    }
    & & \\ \\

    \theutterance \stepcounter{utterance}  
    & & \multicolumn{4}{p{0.6\linewidth}}{
        \cellcolor[rgb]{0.9,0.9,0.9}{
            \makecell[{{p{\linewidth}}}]{
                \texttt{\tiny{[GM$\rangle$P2]}}
                \texttt{Tagad pieņem lēmumu. Kā tu domā: vai tavs režģis un otrs aprakstītais režģis ir vienādi?  Raksti "LĒMUMS: vienāds režģis", ja domā, ka tie ir viens un tas pats režģis, vai "LĒMUMS: atšķirīgi režģi", ja domā, ka tie ir atšķirīgi režģi. Nepievieno neko citu.} \\
            }
        }
    }
    & \\ \\

    \theutterance \stepcounter{utterance}  
    & & & \multicolumn{4}{p{0.6\linewidth}}{
        \cellcolor[rgb]{0.9,0.9,0.9}{
            \makecell[{{p{\linewidth}}}]{
                \texttt{\tiny{[GM$\langle$P2]}}
                \texttt{LĒMUMS: atšķirīgi režģi} \\
            }
        }
    }
     \\ \\

    \theutterance \stepcounter{utterance}  
    & & & \multicolumn{2}{p{0.3\linewidth}}{
        \cellcolor[rgb]{0.9,0.9,0.9}{
            \makecell[{{p{\linewidth}}}]{
                \texttt{\tiny{[GM$|$GM]}}
                \texttt{continue} \\
            }
        }
    }
    & & \\ \\

    \theutterance \stepcounter{utterance}  
    & & & \multicolumn{2}{p{0.3\linewidth}}{
        \cellcolor[rgb]{0.9,0.9,0.9}{
            \makecell[{{p{\linewidth}}}]{
                \texttt{\tiny{[GM$|$GM]}}
                \texttt{success} \\
            }
        }
    }
    & & \\ \\

    \theutterance \stepcounter{utterance}  
    & & \multicolumn{4}{p{0.6\linewidth}}{
        \cellcolor[rgb]{0.9,0.9,0.9}{
            \makecell[{{p{\linewidth}}}]{
                \texttt{\tiny{[P1$\langle$GM]}}
                \texttt{ATBILDE: Nē, pirmajā un piektajā rindā nav X pirmajā un piektajā kolonnā; tās ir tukšas (▢).} \\
\texttt{Tagad pieņem lēmumu. Kā tu domā: vai tavs režģis un otrs aprakstītais režģis ir vienādi?  Raksti "LĒMUMS: vienāds režģis", ja domā, ka tie ir viens un tas pats režģis, vai "LĒMUMS: atšķirīgi režģi", ja domā, ka tie ir atšķirīgi režģi. Nepievieno neko citu.} \\
            }
        }
    }
    & \\ \\

    \theutterance \stepcounter{utterance}  
    & \multicolumn{4}{p{0.6\linewidth}}{
        \cellcolor[rgb]{0.9,0.9,0.9}{
            \makecell[{{p{\linewidth}}}]{
                \texttt{\tiny{[P1$\rangle$GM]}}
                \texttt{LĒMUMS: atšķirīgi režģi} \\
            }
        }
    }
    & & \\ \\

    \theutterance \stepcounter{utterance}  
    & & & \multicolumn{2}{p{0.3\linewidth}}{
        \cellcolor[rgb]{0.9,0.9,0.9}{
            \makecell[{{p{\linewidth}}}]{
                \texttt{\tiny{[GM$|$GM]}}
                \texttt{continue} \\
            }
        }
    }
    & & \\ \\

    \theutterance \stepcounter{utterance}  
    & & & \multicolumn{2}{p{0.3\linewidth}}{
        \cellcolor[rgb]{0.9,0.9,0.9}{
            \makecell[{{p{\linewidth}}}]{
                \texttt{\tiny{[GM$|$GM]}}
                \texttt{success} \\
            }
        }
    }
    & & \\ \\

\end{supertabular}
}

\end{document}
